\documentclass[11pt]{article}
\usepackage[margin=1in]{geometry}
\usepackage[T1]{fontenc}
\usepackage{amsmath,amssymb,amsthm}
\usepackage[hidelinks]{hyperref}
\newtheorem{proposition}{Proposition}
\newcommand{\E}{\mathbb E}
\newcommand{\K}{\mathbb K}
\title{Further outward bridges: Stein operators, apolar Jordan blocks and Hessian duality}
\author{Independent internal research note}
\date{16 November 2026}
\begin{document}
\maketitle
\noindent\textbf{Status.}
These are transfer arguments and research leads. They inherit the assurance
limits of their named repository inputs. No external conjecture, novelty or
priority is claimed solved, no new theorem is promoted to the registry,
and no agents were messaged.

\section{Gaussian certificates become Stein identities}

For a real polynomial $P(X)$ of independent real standard Gaussians, let
$Y=P(X)$ and let $F(t)=\E[e^{tY}]$ denote a \emph{formal} moment series.
Suppose an original Gaussian integration-by-parts certificate gives
\[
       L(t,\partial_t)F=0,\qquad
       L=\sum_{i,j}a_{ij}t^i\partial_t^j.
\]
The anti-isomorphism $t\mapsto\partial_y$, $\partial_t\mapsto y$ sends
this normal-ordered expression to
\[
       S=\sum_{i,j}a_{ij}y^j\partial_y^i.
\]
Coefficientwise, $\E[Sf(Y)]=0$ for every polynomial test $f$. If the
retained original IBP certificate is transported term by term by the
chain rule, the identity also holds for smooth test functions with the
required polynomial-growth and Gaussian integrability conditions. That
extension is a property of the IBP identity, not a consequence merely
of having checked every polynomial moment.

This is exactly the algebraic setting of
\href{https://arxiv.org/abs/1912.04605}{Azmoodeh--Gasbarra--Gaunt,
On algebraic Stein operators for Gaussian polynomials}, and their
\href{https://arxiv.org/abs/2212.07321}{2025 revision on polynomial Stein
operators and Weyl algebras}. Operator existence and bounded-order
searches are already studied there. Our plausible contribution is
replayable original identities and family specialization, not a new
claim of generic existence.

For an explicit control, take $Y=X^3$ and $\theta=y\partial_y$. Then
\[
 S=(3\theta+1)(3\theta+3)(3\theta+5)-y^2
   =27y^3\partial_y^3+162y^2\partial_y^2+
     177y\partial_y+15-y^2.                                  \tag{1}
\]
Indeed Gaussian IBP gives
$\E[X^{2r}\psi]=\E[X^{2r-2}(X\partial_X+2r-1)\psi]$.
Applying it successively with $r=3,2,1$ and $\psi(X)=f(X^3)$ proves
$\E[Sf(X^3)]=0$ for admissible smooth $f$. Equivalently,
\[
 \mu_{n+2}=(3n+1)(3n+3)(3n+5)\mu_n,\quad
 \mu_n=\E[X^{3n}].
\]
Its formal time operator is
$27t^3\partial_t^3+(162t^2-1)\partial_t^2+
177t\partial_t+15$.

A Stein identity need not \emph{characterize} a distribution. Even all
moments need not do so: odd Gaussian powers of degree at least three
are moment-indeterminate. The latter fact is classical, summarized in
\href{https://doi.org/10.1186/s40488-017-0059-2}{Lin's moment-problem survey}.
Our finite formal-jet selectors do not determine an arbitrary probability
law. Characterization requires an admissible analytic solution class,
tail and singular-point conditions, or other distributional restrictions.
The issue is explicitly studied in
\href{https://alea.impa.br/articles/v20/20-06.pdf}{An asymptotic approach to
proving sufficiency of Stein characterisations}.
This is a precise future bridge rather than an already proved analytic
classification.

\section{Apolar algebras: a rational largest-Jordan-block envelope}

Let $\K$ have characteristic zero, and choose a nonconstant polynomial
$P\in\K[x_1,\ldots,x_n]$. Fix a nonzero formal symbol
$\lambda\in(\xi)\K[[\xi_1,\ldots,\xi_n]]$.
For $j\ge0$ form the finite local inverse-system algebra
\[
 A_j=\K[[\xi]]/\operatorname{Ann}(P^j),\qquad
 \operatorname{Ann}(P^j)=\{a:a(\partial_x)P^j=0\}.
\]
It is the usual Artinian Gorenstein apolar algebra, with the possibly
inhomogeneous polynomial $P^j$ as inverse-system generator.
Multiplication by the image of $\lambda$ is nilpotent.
Let $\nu_j$ be its nilpotence index, equivalently its largest Jordan-block
size. This equality holds because multiplication is zero in the $k$th
power exactly when $\lambda^k\cdot1=0$.

The repository's first-survival theorem controls
\[
 \tau_k=\min\{j\ge0:\lambda(\partial)^kP^j\ne0\}.
\]
It states that either $\tau_N=\infty$ for some $N$, or all $\tau_k$
are finite and $\tau_k=\rho k+O(1)$ for a positive rational $\rho$.

\begin{proposition}
The largest-block envelope $H(j)=\max_{0\le q\le j}\nu_q$ is either
eventually constant, or
\[
                       H(j)=j/\rho+O(1)
\]
with positive rational slope. This holds for arbitrary formal $\lambda$
under the retained first-survival theorem.
\end{proposition}
\begin{proof}
Since a nonzero $k$th power of a nilpotent operator implies its previous
powers are nonzero,
\[
 H(j)=1+\max\bigl(\{0\}\cup\{k\ge1:\tau_k\le j\}\bigr).        \tag{2}
\]
If $\tau_N=\infty$, all block sizes are bounded by $N$ and their
nondecreasing integer envelope stabilizes. Otherwise
$\rho k-B\le\tau_k\le\rho k+B$ for large $k$.
Inverting these two inequalities in (2) gives the assertion.
\end{proof}

The envelope is essential: this argument does not prove the same
asymptotic for each individual $\nu_j$, much less its entire Jordan type.

For a direct example, $P=x^N+y^M$ and $\lambda=\xi\eta$ give
\[
 \tau_k=\lceil k/N\rceil+\lceil k/M\rceil,\qquad
 \nu_j=1+\max_{0\le a\le j}\min(Na,M(j-a)).
\]
The envelope slope is $NM/(N+M)$.

This connects our formal results to inverse systems, Jordan types and
Lefschetz questions. A fresh specific target is Conjecture 6.4 of
\href{https://arxiv.org/abs/2610.34517}{Buczynski--Teitler--Woo,
Generators and Lefschetz properties of general hyperplane arrangement
apolar algebras} (25 October 2026): strong Lefschetz for combinatorially
general arrangements in the stated range $n\le d\le2n-2$.
That conjecture concerns ranks of all relevant multiplication maps.
Largest-block information alone does not prove those ranks, and for a
linear $\lambda$ the largest block is elementary directional-degree data.
Thus a transfer to strong Lefschetz requires genuinely new rank control;
no Waring-rank or Lefschetz conjecture follows from this proposition.

\section{An existing geometric theorem gives polynomial Legendre duality}

The repository's Ricci-flat HC4 theorem, which uses the retained Tang
classification, asserts: for $h\in\mathbb C[x_1,\ldots,x_4]$ with constant
nonzero Hessian determinant, vanishing Levi-Civita Ricci tensor of its
\emph{complex bilinear} Hessian metric implies that
$\Phi=\nabla h$ is a polynomial automorphism.

\begin{proposition}
Under these hypotheses, $h$ has a global polynomial Legendre dual
\[
 h^*(p)=p\cdot\Phi^{-1}(p)-h(\Phi^{-1}(p)).
\]
It has constant nonzero Hessian determinant, its gradient is
$\Phi^{-1}$, and its complex bilinear Hessian metric is again Ricci-flat.
\end{proposition}
\begin{proof}
Differentiate the displayed expression to obtain
\[\nabla h^*(p)=\Phi^{-1}(p),\qquad
\operatorname{Hess}h^*(p)=
(\operatorname{Hess}h(\Phi^{-1}(p)))^{-1}.\]
The gradient map pulls this latter metric back to
$\operatorname{Hess}h$, so they are isometric.
\end{proof}

This is a direct outward consequence for the polynomial Monge--Ampere
equation $\det\operatorname{Hess}h=c$ and algebraic Legendre duality.
It does not classify every such polynomial. Complex bilinear
Levi-Civita Ricci-flatness must not be replaced by Hermitian Ricci
flatness, scalar curvature zero, or the automatically vanishing
expression $-\partial_i\partial_j\log\det\operatorname{Hess}h$.
The separate repository examples already show that polynomial invertibility
does not force the Hessian metric to be flat or geodesically complete.

\section{Formal finite-dimensional physics and exact integration}

The full Gaussian IBP source is the Schwinger--Dyson system of the
formal finite-dimensional action
$S(X)=\sum_iX_i^2/2-tP(X)$. Its amplitudes $Q$ are observable insertions.
For each fixed observable,
\[
       Z_Q(t)=\E[Qe^{tP}]
             =\sum_{m\ge0}\E[QP^m]t^m/m!.
\]
Thus Gaussian all-marker safety says that every fixed polynomial
observable has a terminating formal interaction expansion. The field-safe
ideal and its coefficient-ring normalization now classify that condition
on specified finite coefficient families, subject to their implementation
and assurance limits.

The retained quartic counterexample has $Z_1(t)=1$ but, for its original
marker $Q=Z$, $Z_Z(t)=t/(1-t)$. Hence cancellation of every positive
vacuum coefficient does not imply termination with insertions.
This is an interpretation of a known witness, not a new physical model.

The IBP/Schwinger--Dyson identity is standard, for example in
\href{https://www.damtp.cam.ac.uk/user/dbs26/AQFT/chaps1\%2B2.pdf}{Skinner's
zero-dimensional field-theory notes}. Related D-module methods for genuine
Feynman integrals are developed by
\href{https://doi.org/10.1007/s11005-024-01835-7}{Henn--Pratt--Sattelberger--Zoia}.
Passing to propagator denominators, singular integration cycles,
renormalization or infinitely many field variables requires further work.
No continuum field-theory result is asserted.

\section{Singularities: a plausible transfer with an important status correction}

Our log-discrepancy, boundary-pole and Bernstein machinery belongs to the
same mathematical setting as monodromy conjectures for hypersurface
singularities. It may be useful in selected classes beyond arrangements.
However, the hyperplane-arrangement case is not a clean untouched target:
May 2026 preprints by
\href{https://arxiv.org/abs/2605.25335}{Wu} and
\href{https://arxiv.org/abs/2605.25314}{Davis--Yang} claim the strong
monodromy statement there. Their proofs were not audited here.
An actual new application would have to specify a hypersurface class and
prove the missing link from valuation data to the required Bernstein root.

\section{Priority assessment}

The closest useful new application is certified Stein identities and
their specialization behavior. Distribution characterization is a separate
analytic target. The apolar envelope and Ricci-flat polynomial Legendre
duality are precise consequences, rather than speculative analogies.
The Schwinger--Dyson reading is exact for formal finite-dimensional
integrals. Larger Feynman, strong Lefschetz and general monodromy problems
still require substantial new bridges.
\end{document}
